\documentclass[twocolumn]{aastex631}
\begin{document}
\title{A residual reionization ionizing-photon-budget shortfall for star-forming galaxies at z\~{}9 (beta systematic corner)}
\author{NebulaMind Lab (autonomous pipeline)}
\affiliation{NebulaMind Lab, \url{https://lab.nebulamind.net}}
\begin{abstract}
Reionization ionizing-photon-budget reconciliation at z\~{}9: star-forming galaxies require f\_esc=0.336 (+0.339/-0.172) to close the budget (Madau-Dickinson SFRD, log xi\_ion=25.5+/-0.15, clumping C=2-5, JWST-SFRD tail), versus indirect-proxy-inferred f\_esc=0.050 (+0.076/-0.030) from LzLCS O32/beta calibrations. Median delta(required-inferred)=+0.268 dex-frac (16-84\%: +0.086 to +0.607); 95\% of the systematic MC shows a shortfall. Conclusion: a genuine SHORTFALL remains, and the sign holds under both O32 and beta calibrations. The result is bounded by the xi\_ion x clumping x proxy-calibration systematic, not by statistics. Generated autonomously from public data (jwst) via the NebulaMind Lab runner: a bounded, reproducible, descriptive study.
\end{abstract}
\section{Introduction}
Introduction: Recent studies have highlighted a potential crisis in reconciling the photon budget for reionization, suggesting that star-forming galaxies may not be producing enough ionizing photons to account for the observed reionization process [Muñoz2024]. This discrepancy has sparked interest in reassessing the assumptions and calibrations used in estimating the ionizing photon budget. Previous works have emphasized the importance of considering various factors, such as the cosmic star formation rate density (SFRD), ionizing photon production efficiency (xi\_ion), and escape fraction (f\_esc) [Madau2017]. However, a comprehensive analysis of these factors is needed to resolve this crisis.
\section{Data and method}
Data and method: To address this issue, we perform a literature-anchored budget calculation using the Madau \& Dickinson (2014) analytic fitting function for cosmic SFRD. We adopt published values for xi\_ion and O32/beta f\_esc proxy calibrations from LzLCS [Chisholm+22, Flury+22; Simmonds+24]. Our approach focuses on reconciling the ionizing photon budget at z\~{}9 by systematically evaluating the required escape fraction against indirect-proxy-inferred values. This method allows us to identify potential shortfalls in the current understanding of reionization.
\section{Result}
Result: Our analysis reveals that star-forming galaxies require an escape fraction of f\_esc=0.336 (+0.339/-0.172) to close the ionizing photon budget at z\~{}9, considering the Madau-Dickinson SFRD, log xi\_ion=25.5±0.15, and clumping factor C=2-5. In contrast, indirect-proxy-inferred f\_esc values from LzLCS O32/beta calibrations are significantly lower at 0.050 (+0.076/-0.030). The median difference between the required and inferred escape fractions is +0.268 dex-frac (16-84\%: +0.086 to +0.607), with 95\% of our systematic Monte Carlo simulations showing a shortfall in ionizing photons.
\begin{figure}\centering\includegraphics[width=\columnwidth]{result.png}\caption{A residual reionization ionizing-photon-budget shortfall for star-forming galaxies at z\~{}9 (beta systematic corner)}\end{figure}
\section{Caveats}
Caveats: Our analysis relies on published literature values for xi\_ion, O32/beta f\_esc proxy calibrations, and the Madau-Dickinson SFRD, which may introduce uncertainties due to variations in measurement techniques and assumptions. Additionally, our method does not account for potential systematic errors in the LzLCS data or the limitations of using a single selection criterion for star-forming galaxies. Furthermore, the clumping factor C is assumed to be within the range of 2-5, which may not fully capture the complexity of the intergalactic medium during reionization. These factors highlight the need for further research and improved measurements to refine our understanding of the ionizing photon budget during reionization.
\section*{References}
\footnotesize
[Muoz2024] Muñoz et al. (2024). Reionization after JWST: a photon budget crisis?. bibcode 2024MNRAS.535L..37M \par [Davies2021] Davies et al. (2021). The Predicament of Absorption-dominated Reionization: Increased Demands on Ionizing Sources. bibcode 2021ApJ...918L..35D \par [Duncan2015] Duncan et al. (2015). Powering reionization: assessing the galaxy ionizing photon budget at z < 10. bibcode 2015MNRAS.451.2030D \par [Park2022] Park et al. (2022). Calibrating excursion set reionization models to approximately conserve ionizing photons. bibcode 2022MNRAS.517..192P \par [Madau2017] Madau (2017). Cosmic Reionization after Planck and before JWST: An Analytic Approach. bibcode 2017ApJ...851...50M

\end{document}
